
Reinforcement Learning from Human Feedback (RLHF) and Direct Preference Optimization (DPO) represent two approaches for aligning language models with human preferences. These methods differ in their optimization paths: RLHF trains an explicit reward model that learns a continuous approximation of human preferences through Bradley-Terry loss, while DPO directly optimizes the policy from preference pairs using a closed-form objective without an intermediate reward signal.

This mechanistic divergence raises a question: do these different training dynamics produce measurably different alignment signatures? The theoretical motivation is that RLHF's Bradley-Terry loss learns continuous reward values that interpolate between discrete preference labels, creating smooth optimization landscapes. DPO's closed-form objective directly encodes preference rankings into policy updates, potentially preserving sharper preference boundaries. If optimization landscape geometry influences convergence behavior, these methods might converge to different stable configurations that manifest as differential performance patterns across alignment benchmarks.

This hypothesis was tested under controlled conditions: same base model (Llama-2-7B), same preference data (Anthropic HH-RLHF), same compute budget. A five-hypothesis experimental design was used to isolate the mechanistic differences and trace their effects through to benchmark-level evaluation.

The findings reveal an unexpected gap in the causal chain. RLHF and DPO exhibit genuinely different training dynamics: RLHF produces smooth, continuous reward predictions spanning a 0.83-unit range, while DPO shows 65\% higher margin variance (sharpness ratio = 1.65). These mechanistic differences are measurable.

However, the predicted downstream effects did not follow. Models trained with RLHF and DPO did not cluster into distinct behavioral attractors---cross-method similarity exceeded within-method similarity (clustering gap = -0.016). Standard alignment benchmarks did not reveal large differential profiles (max |Cohen's d| = 0.194, below the 0.3 threshold for a medium effect).

\textbf{Contributions:}
\begin{enumerate}
\item Empirical verification of RLHF reward model smoothness versus DPO preference boundary sharpness on identical data and model
\item Controlled test of the behavioral attractor hypothesis at 7B scale, which was not confirmed
\item Evidence that standard alignment benchmarks may lack sensitivity to method-specific differences
\item Validation that TruthfulQA and HHH benchmarks measure independent alignment dimensions (all pairwise |r| < 0.05)
\end{enumerate}

